\documentclass[10pt,twocolumn]{article}

\usepackage[margin=0.72in]{geometry}
\usepackage[T1]{fontenc}
\usepackage[utf8]{inputenc}
\usepackage{url}

\setlength{\columnsep}{0.25in}
\newcommand{\preserved}{\textsc{preserved}}
\newcommand{\partialpreserved}{\textsc{partially preserved}}
\newcommand{\verifiedpreserved}{\textsc{verified preserved}}
\newcommand{\observedpartial}{\textsc{observed partial}}
\newcommand{\altered}{\textsc{altered}}
\newcommand{\lost}{\textsc{lost}}
\newcommand{\confirmedlost}{\textsc{confirmed lost}}
\newcommand{\unresolved}{\textsc{unresolved equivalence}}
\newcommand{\measured}{\textsc{measurement error}}

\title{A11yPreserve: An Uncertainty-Aware Conformance Suite for DOCX-to-PDF Accessibility Preservation}
\author{Anonymous for review}
\date{}

\begin{document}
\maketitle

\begin{abstract}
Accessible source documents can contain author-supplied information such as heading structure, alternative text, language metadata, list semantics, table relationships, and footnote associations. When a document is converted, inspecting only the destination does not establish whether those source properties survived. We present A11yPreserve, an uncertainty-aware conformance suite that freezes controlled source contracts, independently checks their DOCX representations, converts identical sources through real pipelines, extracts destination PDF structure through multiple evidence paths, and compares source and destination properties with provenance-traceable evidence. The primary benchmark contains 13 atomic DOCX fixtures and 26 genuine DOCX-to-PDF outputs from LibreOffice and the Google Docs DOCX-to-PDF conversion pipeline. Its evidence-aware results are 11 verified preserved, 3 observed partial, 3 altered, 2 confirmed lost, 6 unresolved equivalence, and 1 measurement error. The two confirmed losses concern inline language and footnote association in the tested Google Docs pipeline. A separate three-document integrated sanity check shows that the protocol remains usable when existing feature contracts coexist. Results are feature-specific and are not used to rank converters globally. The study demonstrates that source-grounded preservation is measurable and that tagged destination output alone does not fully characterize whether known source accessibility information survived conversion.
\end{abstract}

\noindent\textbf{Keywords:} document conversion, accessibility preservation, DOCX, PDF, tagged PDF, differential testing, alternative text, document structure

\section{Introduction}
\label{sec:introduction}

Authors can encode accessibility information directly in source documents: a heading is more than large text, an image may carry an author-written alternative, a run may identify a language change, and a table may identify header relationships. These properties matter because the source document is often exported, uploaded, or converted before it reaches readers. The resulting file may contain some accessibility structure, but that fact alone does not show whether the information intentionally supplied by the author survived the transformation.

This distinction is recognized in accessibility guidance. ATAG 2.0 addresses preservation during content transformations when equivalent output mechanisms exist \cite{atag20}. Revised Section 508 includes a requirement to preserve information provided for accessibility during format conversion to the extent supported by the destination format \cite{section508}. EN 301 549 likewise specifies preservation of accessibility information in transformations \cite{etsi301549}. These requirements motivate a practical question: how can preservation be tested against known source ground truth rather than inferred from the destination alone?

Prior work has studied office-document conversion, PDF-to-HTML accessibility, automated PDF remediation, and destination accessibility evaluation. Roig and Ribera examined office-to-EPUB conversion \cite{roig2015}; SciA11y and related Paper-to-HTML work reconstruct accessible representations from scientific PDFs \cite{wang2021scia11y}; iTagPDF uses source cues to improve PDF tagging \cite{mowar2026}; and recent benchmarks evaluate PDF accessibility or conversion services \cite{kumar2025,daisy2026,docaccessible2026,docaccessible-google}. A11yPreserve complements these efforts by beginning with controlled DOCX documents whose target accessibility properties are explicitly verified, then measuring whether those properties survive two real DOCX-to-PDF pipelines. The contribution is the controlled source-to-destination preservation protocol and empirical benchmark, not a claim that accessible conversion has not been studied.

We present A11yPreserve as a source-contract conformance protocol, benchmark, and empirical study. The workflow is:

\begin{center}
\fbox{Verified DOCX source}\ $\rightarrow$\ 
\fbox{Real converter}\ $\rightarrow$\ 
\fbox{Tagged PDF}\ $\rightarrow$\ 
\fbox{Structure extraction}\ $\rightarrow$\ 
\fbox{Source-grounded comparison}
\end{center}

The primary benchmark contains 13 atomic DOCX fixtures and 26 genuine PDF outputs: 13 from LibreOffice and 13 from the Google Docs DOCX-to-PDF conversion pipeline. All outputs were parseable, non-empty, tagged, and passed the predefined structural output checks. The evidence-aware results are 11 \verifiedpreserved{}, 3 \observedpartial{}, 3 \altered{}, 2 \confirmedlost{}, 6 \unresolved{}, and 1 \measured{}. A separate secondary sanity check evaluates three realistic multi-feature documents using the same contracts and pipelines; it is not merged into the primary denominator. The two confirmed losses concern inline language and footnote association in the Google Docs pipeline. These findings are not a product ranking and do not estimate the prevalence of conversion problems in arbitrary documents. The historical V1 counts remain archived for provenance rather than serving as the primary result summary.

The primary research question is: \emph{Does accessibility information present in a source document survive document conversion?} The paper makes three methodological contributions and one bounded empirical application:

\begin{enumerate}
\item versioned, feature-specific source contracts for controlled DOCX accessibility properties;
\item an independently checked source oracle and destination-evidence protocol that separates verified equivalence, observed difference, unresolved equivalence, confirmed absence, and measurement failure;
\item a reusable conformance-suite design with calibrated differential comparison, evidence certificates, and provenance; and
\item an empirical application to 26 real outputs from two DOCX-to-PDF pipelines, plus a separate three-document integrated sanity check, reporting preservation behavior only for the tested fixtures, documents, and pipeline conditions.
\end{enumerate}

The paper deliberately makes no universal claim about accessibility loss during conversion, does not rank either converter globally, and does not treat the results as real-world prevalence. It asks a simpler question: when accessibility information is known to exist before conversion, did that information survive?



\section{Background and Related Work}
\label{sec:background}

\subsection{Preservation as an accessibility requirement}

Accessibility information is often authored in one representation and transformed into another. ATAG 2.0 explicitly addresses restructuring and recoding transformations: when equivalent mechanisms exist in the output technology, accessibility information should be preserved or the author should be warned or prompted to check it \cite{atag20}. Its text-alternative criterion is especially close to the F02 contract. The U.S. Revised Section 508 standards similarly state that authoring tools should preserve information required for accessibility during format conversion to the extent supported by the destination format \cite{section508}. EN 301 549 contains a corresponding preservation-of-accessibility-information requirement for transformations \cite{etsi301549}. These requirements motivate testing preservation, but they do not predict the outcomes of the converters evaluated here.

\subsection{Accessibility-aware document conversion}

Roig and Ribera evaluated office-document conversion to EPUB and reported accessibility-preservation problems in that setting \cite{roig2015}. Their work is an important precedent for studying accessibility effects of conversion. A11yPreserve differs in experimental scope and instrumentation: it uses frozen native DOCX source contracts, real DOCX-to-PDF pipelines, destination structure extraction, and provenance-traceable fixture-level differential records. The distinction is one of scope and method, not a claim that prior work ignored preservation.

\subsection{Accessible representation generation and remediation}

SciA11y converts scientific papers from PDF to accessible HTML and combines extraction with navigation-oriented rendering for blind and low-vision readers \cite{wang2021scia11y}. This work primarily creates a more accessible destination representation. iTagPDF uses source-document cues and visual PDF analysis to automate tagging, reading order, and content-specific metadata such as alt text and formula structure \cite{mowar2026}. Its source-aware remediation objective is closely related to A11yPreserve's use of source information, but the experimental question differs: iTagPDF generates or improves PDF metadata, whereas A11yPreserve measures whether author-provided source information survived an existing conversion pipeline.

\subsection{PDF conversion and accessibility evaluation benchmarks}

Recent work also benchmarks destination conversion or evaluation. Kumar, Padath, and Wang introduce a benchmark for assessing automated and LLM-based PDF accessibility evaluation across criteria including alternative text, reading order, tagging, tables, and hyperlinks \cite{kumar2025}. DAISY's 2026 benchmark compares seven services that convert PDFs into accessible documents against a shared quality framework \cite{daisy2026}. DocAccessible's public research program reports source-grounded PDF-to-HTML conversion and fidelity benchmarking with frozen evidence \cite{docaccessible2026}. Its current Google Docs guidance also emphasizes that tagged export requires post-export checks of language, title, table headers, and reading order, and does not itself assert WCAG or PDF/UA conformance \cite{docaccessible-google}. These are important adjacent efforts. A11yPreserve's narrower contribution is to start from controlled DOCX accessibility ground truth and compare the same source semantics across real DOCX-to-PDF pipelines, while separating preservation fidelity from destination-only accessibility evaluation.

Vendor documentation provides useful format context as well. Microsoft's current Word documentation maps source constructs such as language spans, footnotes, equations, captions, hyperlinks, and decorative content to PDF structure, including \texttt{/Lang}, \texttt{/Note}, \texttt{/Formula}, \texttt{/Caption}, and artifacts \cite{microsoft-word-pdf}. We use such documentation to define representability and scope, not as comparative evidence about the tested pipelines.

The resulting contribution boundary is therefore modest: not the invention of accessibility-conversion benchmarking, but a source-grounded DOCX-to-PDF preservation conformance design with explicit contracts, independently checked source hashes, calibrated feature-level evidence statuses, and a reproducible evidence chain. The contribution is the controlled source-to-destination preservation protocol and empirical benchmark; it is not a claim that prior work lacked accessible-conversion studies.



\section{A11yPreserve Method}
\label{sec:method}

A11yPreserve treats accessibility preservation as a source-to-destination comparison problem. The method begins with a document whose target accessibility properties are known, runs that exact source through a named conversion pipeline, extracts accessibility-relevant destination structure, and classifies the correspondence. This framing is intentionally narrower than asking whether a PDF is accessible in every respect. A destination-only checker can report tags or metadata that are useful in their own right, but it cannot by itself establish whether an author-provided source property was retained.

\subsection{Source fixtures and ground truth}

The frozen corpus contains 13 small DOCX fixtures. The fixtures cover heading hierarchy, image alternative text, list semantics, table/header structure, document language, inline language changes, decorative image state, hyperlinks, document title and metadata, multi-level table headers, footnote association, equation/math semantics, and figure/caption association. They were created by controlled generation scripts using \texttt{python-docx} together with explicit WordprocessingML constructs where required. Each fixture isolates a target property and avoids unnecessary layout complexity.

Each fixture has a machine-readable preservation contract under \texttt{contracts/}. The frozen source-grounding chain used three artifacts: the source manifest, raw OOXML inspection, and the source extractor's accessibility intermediate representation (ASIR). The strengthening pass added a fourth, separate audit path: an independent source-oracle implementation performs fresh ZIP/XML assertions without importing ASIR, fixture-generation, manifest-generation, or comparator code; its per-fixture certificates are retained under \texttt{evidence/source/}. OOXML checks inspect the relevant document parts, including document XML, run properties, numbering, relationships, and the footnotes part where applicable. The ASIR remains an implementation-grounded oracle, not an independent implementation or independent human expert validation: it is part of the same analysis toolchain. Agreement among the frozen manifest, raw OOXML, independent oracle, and ASIR provides a stronger cross-check against extractor-only errors, but does not remove the limitation that no external expert panel validated every contract. This establishes encoded source ground truth; it does not claim validation by disabled users. The corpus manifest records file size, SHA-256, feature family, and verification status.
The independent OOXML oracle agreed with all 13 frozen source contracts.

\subsection{Conversion pipelines}

The first condition used LibreOffice 26.2.6.3, build \texttt{8221e31b3ac356a1623c672912a3d2b492f7e3d1}, on Windows \texttt{10.0.26200.0}. Conversion used native headless \texttt{soffice.com writer\_pdf\_Export} with \texttt{UseTaggedPDF=true}, \texttt{PDFUACompliance=true}, and \texttt{SelectPdfVersion=1}. The second condition used an authenticated Google Docs web workflow on the same Windows environment with Chrome 153.0.8010.53: upload the exact frozen DOCX and choose \textit{File} $\rightarrow$ \textit{Download} $\rightarrow$ \textit{PDF Document (.pdf)}. The observed Google PDF producer was \texttt{Skia/PDF m156}. No print-to-PDF path, screenshot, manual content repair, or post-export accessibility repair was used. The official LibreOffice documentation describes its PDF/UA export option \cite{libreoffice-pdfua}; Google's download documentation describes the selected-file-type download workflow \cite{google-docs-download}.

The Google condition is described throughout as the \emph{Google Docs DOCX-to-PDF conversion pipeline}. It includes DOCX import, Google's internal document representation, and PDF generation. The backend version cannot be fully pinned, so the experiment date matters and a rerun may change after a service update; the workflow, browser, operating system, observed producer, and environment were recorded. The experiment does not isolate which internal stage caused an observed change. Microsoft Word was not included in the denominator because it was not a completed primary pipeline.

\subsection{Destination extraction and validation}

All 26 outputs were checked for existence, non-zero size, parseability, expected page count, structural tagging, and a PDF structure tree. The extractor inspected PDF structure elements and roles relevant to the fixture contracts: headings, figures and \texttt{/Alt}, list containers and \texttt{/ListNumbering}, tables and header cells, \texttt{/Lang}, link annotations, metadata, \texttt{/Note}, \texttt{/Formula}, and \texttt{/Caption}. Evidence was retained as structured JSON and raw object observations. Poppler \texttt{pdftotext} was used as an independent secondary check when marked-content text was not recoverable from the structural parser.

The extractor is deliberately not treated as infallible. For example, the LibreOffice F11 output contains a \texttt{/Note} structure, but the available extraction evidence does not establish the note-body association. That case is retained as \measured{} rather than converted into a positive or negative preservation claim. For F06, F11, and the difficult partial cases, an independent triangulation path uses PyMuPDF page/text/link extraction and serialized-object inspection in addition to pypdf structure traversal. This reduces, but does not eliminate, parser uncertainty.

\subsection{Differential comparison}

The \texttt{a11ydiff} implementation loads the source ASIR and destination extraction for each fixture/pipeline pair. It records the source value, destination value, destination representation, evidence paths, source/output hashes, confidence, and a classification. The feature contracts and classification rules were fixed before the final conversion results were reviewed. A separate synthetic calibration set exercises positive, negative, and alternate-equivalent destination controls; its diagnostic counts are not population estimates and reveal conservative false-negative/indeterminate behavior for some alternate representations. The implementation also records a secondary destination-accessibility axis. The final evidence-aware result labels are defined below.

\begin{description}
\item[\verifiedpreserved] The destination expresses the contracted property exactly or through a verified equivalent representation, with independent source and destination evidence.
\item[\observedpartial] Some contract components survive and at least one required representable component is directly demonstrated to differ or be absent.
\item[\altered] The feature remains present, but an author-relevant value changes, such as alt-text wording, locale specificity, or title metadata.
\item[\confirmedlost] The source property is valid and representable in the destination, but the required destination representation is absent and independent checks support that absence.
\item[\unresolved] Relevant destination structure exists, but available evidence cannot establish whether the source contract survived through an alternate PDF representation; absence has not been independently demonstrated.
\item[\measured] The output is valid and relevant structure is present, but the available evidence cannot resolve the contract reliably.
\end{description}

The historical V1 result set used five categories: PRESERVED, PARTIALLY\_PRESERVED, ALTERED, LOST, and MEASUREMENT\_ERROR. The primary manuscript result model is the six-way evidence-aware taxonomy above; V1 remains archived for provenance. The source and destination files did not change.

\subsection{Independent verification and freezing}

Non-preserved cases were reviewed against the manifest, raw OOXML, independent source certificates, source ASIR, output validation, primary PDF extraction, full structure traversal, and independent parser/object checks where appropriate. The two \confirmedlost{} cases were retained only after two destination evidence paths corroborated absence of the required machine-verifiable representation; visible text survival is not treated as preservation of the associated structure. Source fixtures and PDF outputs were SHA-256 frozen. Phase 6 independently recomputed all 13 source hashes and 26 output hashes, normalized the primary and \texttt{a11ydiff} classification records, and reconstructed all 26 rows without discrepancy. The V2 evidence-aware result file is derived from the frozen V1 records and does not change fixtures, PDFs, manifests, hashes, or denominators.



\section{Experimental Design}
\label{sec:experiment}

\subsection{Research questions}

The experiment addresses three questions:

\begin{enumerate}
\item \textbf{RQ1:} To what extent do the tested DOCX-to-PDF conversion pipelines preserve accessibility information present in controlled source documents?
\item \textbf{RQ2:} Which tested accessibility properties are verified as preserved, observed as changed or lost, or unresolved under the available cross-format evidence?
\item \textbf{RQ3:} How do preservation outcomes differ between the tested LibreOffice and Google Docs conversion pipelines?
\end{enumerate}

Validator comparison is not a fourth core question. The project did not obtain a systematic PAC, Acrobat Accessibility Checker, or veraPDF dataset for all 26 outputs, so destination-only validator comparison is future work rather than a central result.

\subsection{Benchmark composition}

Table~\ref{tab:fixtures} summarizes the 13 atomic fixture contracts. Each fixture contributes one feature-level outcome per pipeline, yielding 13 $\times$ 2 = 26 cases. The fixtures are intentionally controlled rather than randomly sampled. F13 was reserved for reading order but deferred because reliable automated source-to-PDF equivalence could not be established; the frozen corpus therefore uses F01--F12 and F14 without renumbering. The table's contract descriptions state what is compared; they do not imply that the feature families have equal complexity or equal user impact.

\begin{table*}[t]
\centering\small
\caption{Frozen fixture contracts.}
\label{tab:fixtures}
\begin{tabular}{p{0.14\textwidth}p{0.17\textwidth}p{0.56\textwidth}}
\hline
Fixture & Property & Contract focus \\
\hline
F01 & Headings & Heading levels and order \\
F02 & Alternative text & Author-written image description \\
F03 & Lists & Ordered/unordered type, nesting, and list structure \\
F04 & Table & Table, row/cell, and header structure \\
F05 & Document language & Primary document language value \\
F06 & Inline language & Run-level language override within an en-US document \\
F07 & Decorative image & Meaningful versus explicitly decorative image state \\
F08 & Hyperlinks & Link existence, visible text, and URI destination \\
F09 & Document title & Visible title and document metadata identity \\
F10 & Complex table & Multi-level headers, spans, and associations \\
F11 & Footnotes & Note body and reference-to-note association \\
F12 & Equation & Native equation representation and visible expression \\
F14 & Caption & Figure/caption text and association \\
\hline
\end{tabular}
\end{table*}

\subsection{Converter conditions and provenance}

The LibreOffice workflow was deterministic and local. The Google Docs workflow was authenticated and browser-mediated, so its service backend was not independently version-pinned. The experiment date, browser version, observed PDF producer, source hashes, output hashes, file sizes, and validation records were retained in the automation log. All outputs were one-page PDFs in this corpus and passed the output-validation checks.

The experiment's primary unit is one fixture/converter pair. The 13 fixtures are controlled atomic cases, not a representative sample; features have unequal complexity and are not repeated instances of a population. Counts and feature-level matrices are therefore primary, while percentages are not used as headline results. No inferential statistics, significance tests, prevalence estimates, or weighted overall converter score are justified. A separate secondary sanity check uses three realistic multi-feature documents and the same existing contracts; its property-level observations are not merged into the primary denominator.

\subsection{Outcome accounting}

The primary evidence-aware set contains 11 \verifiedpreserved{}, 3 \observedpartial{}, 3 \altered{}, 2 \confirmedlost{}, 6 \unresolved{}, and 1 \measured{} case. Thus 26 cases were attempted and successfully converted, with one case remaining a measurement error. The historical V1 set contained 11 PRESERVED, 9 PARTIALLY\_PRESERVED, 3 ALTERED, 2 LOST, and 1 MEASUREMENT\_ERROR; the V1 record is retained in the interpretation changelog and is not silently overwritten. No percentages are used because the controlled feature units have unequal complexity and do not support prevalence interpretation.

\subsection{Integrated-document sanity check}

The secondary sanity check evaluated three realistic two-page documents: I01, a community access report; I02, an educational handout; and I03, a policy note. They used only existing contract families, and an independent raw-OOXML oracle passed 21/21 integrated source assertions. The same two pipelines produced six tagged PDFs, and the protocol generated 42 property--pipeline observations without requiring a document-level score. The secondary result distribution was 16 VERIFIED\_PRESERVED, 6 OBSERVED\_PARTIAL, 5 ALTERED, 3 CONFIRMED\_LOST, 11 UNRESOLVED\_EQUIVALENCE, and 1 MEASUREMENT\_ERROR. These observations demonstrate protocol applicability when properties coexist; they are not a prevalence estimate and are excluded from the primary 26-case denominator.



\section{Results}
\label{sec:results}

\subsection{Overall accounting}

All 26 conversion attempts produced genuine, parseable, non-empty, structurally tagged PDFs that passed the predefined structural output checks. The final evidence-aware result counts are 11 \verifiedpreserved{}, 3 \observedpartial{}, 3 \altered{}, 2 \confirmedlost{}, 6 \unresolved{}, and 1 \measured{}. The historical V1 distribution was 11 PRESERVED, 9 PARTIALLY\_PRESERVED, 3 ALTERED, 2 LOST, and 1 MEASUREMENT\_ERROR; it remains archived for provenance in the interpretation changelog. The V2 model makes uncertainty primary rather than folding unresolved equivalence into observed partial preservation.

\begin{figure*}[t]
\centering
\fbox{\begin{minipage}{0.91\textwidth}
\centering
\textbf{Verified DOCX source}\quad$\longrightarrow$\quad\textbf{Real conversion pipeline}\quad$\longrightarrow$\quad\textbf{Tagged PDF}\quad$\longrightarrow$\quad\textbf{PDF structure extraction}\quad$\longrightarrow$\quad\textbf{\texttt{a11ydiff} comparison}\\[3pt]
Manifest + OOXML + ASIR\hfill Converter/version/settings\hfill StructTreeRoot and roles\hfill Source contract versus destination representation\\[3pt]
\centering\textsc{verified preserved}\quad|\quad\textsc{observed partial}\quad|\quad\textsc{altered}\quad|\quad\textsc{confirmed lost}\quad|\quad\textsc{unresolved}\quad|\quad\textsc{measurement error}
\end{minipage}}
\caption{A11yPreserve's source-grounded measurement workflow. The destination is inspected, but the outcome is defined by comparison with verified source information.}
\label{fig:workflow}
\end{figure*}

\subsection{Feature-level outcomes}

Table~\ref{tab:results} is the complete frozen feature-level result matrix using the final evidence-aware labels. LibreOffice verified preservation for F01, F03, F04, F05, F06, F09, F10, F12, and F14; altered F02; showed observed partial preservation for F07; left F08 unresolved; and produced the F11 measurement error. The Google Docs pipeline verified preservation for F02 and F04, showed observed partial preservation for F01 and F07, left F03, F08, F10, F12, and F14 unresolved, altered F05 and F09, and confirmed loss for F06 and F11. These are feature-specific observations, not a global product ranking.

The V2 model identifies three \emph{observed partial} cases: Google Docs F01 has an extra H1 in the destination heading sequence, and both F07 outputs emit the image pair as Figure elements while directly differing from the source decorative-state contract. Six cases are \emph{unresolved equivalence}: Google Docs F03, F10, F12, and F14, plus both F08 cases. In those cases, relevant structure or visible content exists, but available cross-format evidence does not establish the required type, association, scope, or native semantic representation; absence has not been independently demonstrated. None is a measurement error; the sole measurement error remains F11/LibreOffice.

\begin{table*}[t]
\centering\scriptsize
\caption{Complete primary fixture/converter matrix using the final evidence-aware result labels. ``Partial'' means directly observed component-level difference; ``unresolved'' means relevant structure exists but cross-format equivalence is not established; ``lost'' means independently corroborated absence.}
\label{tab:results}
\begin{tabular}{p{0.19\textwidth}p{0.18\textwidth}p{0.16\textwidth}p{0.35\textwidth}}
\hline
Feature & Pipeline & Final result & Key evidence \\
\hline
F01 headings & LibreOffice & Preserved & levels/order verified \\
F01 headings & Google Docs & Partial & extra H1 changes sequence \\
F02 alt text & LibreOffice & Altered & wording regenerated \\
F02 alt text & Google Docs & Preserved & source string retained \\
F03 lists & LibreOffice & Preserved & L/LI and list type \\
F03 lists & Google Docs & Unresolved & nesting survives; type not established \\
F04 table & LibreOffice & Preserved & table/cell/header roles \\
F04 table & Google Docs & Preserved & table/cell/header roles \\
F05 document language & LibreOffice & Preserved & language value \\
F05 document language & Google Docs & Altered & en-US becomes en \\
F06 inline language & LibreOffice & Preserved & es-MX span \\
F06 inline language & Google Docs & Lost & inline absence corroborated \\
F07 decorative image & LibreOffice & Partial & Figure/Alt differs from decorative contract \\
F07 decorative image & Google Docs & Partial & empty decorative Figure differs \\
F08 hyperlinks & LibreOffice & Unresolved & URI survives; name association unresolved \\
F08 hyperlinks & Google Docs & Unresolved & URI survives; name association unresolved \\
F09 title/metadata & LibreOffice & Preserved & title/metadata \\
F09 title/metadata & Google Docs & Altered & metadata title changes \\
F10 complex table & LibreOffice & Preserved & scope/span/header evidence \\
F10 complex table & Google Docs & Unresolved & associations not established \\
F11 footnotes & LibreOffice & Measurement error & Note exists; association unmeasured \\
F11 footnotes & Google Docs & Lost & visible note; no Note association \\
F12 equation & LibreOffice & Preserved & Formula structure \\
F12 equation & Google Docs & Unresolved & visible expression; math structure unresolved \\
F14 captions & LibreOffice & Preserved & Caption structure \\
F14 captions & Google Docs & Unresolved & caption association unresolved \\
\hline
\end{tabular}
\end{table*}

\begin{figure}[t]
\centering
\scriptsize
\begin{tabular}{lr}
VERIFIED PRESERVED & 11\\[-1pt]
\rule{0.85\linewidth}{1.1ex} & \\
OBSERVED PARTIAL & 3\\[-1pt]
\rule{0.23\linewidth}{1.1ex} & \\
ALTERED & 3\\[-1pt]
\rule{0.23\linewidth}{1.1ex} & \\
CONFIRMED LOST & 2\\[-1pt]
\rule{0.15\linewidth}{1.1ex} & \\
UNRESOLVED & 6\\[-1pt]
\rule{0.46\linewidth}{1.1ex} & \\
MEASUREMENT ERROR & 1\\[-1pt]
\rule{0.08\linewidth}{1.1ex} & \\
\end{tabular}
\caption{Final evidence-aware outcome counts across 26 fixture/pipeline cases. Counts describe controlled cases and are not prevalence estimates or weighted accessibility scores.}
\label{fig:outcomes}
\end{figure}

\subsection{Cross-pipeline differences}

There are 10 fixture rows where the two pipelines have different frozen primary outcomes. The difference is not a single monotonic pattern. For headings, the LibreOffice sequence matched the source while Google Docs added an H1. For alternative text, Google Docs retained the exact author string while LibreOffice prefixed it. For lists, LibreOffice emitted explicit \texttt{/ListNumbering} evidence while Google Docs retained nesting without independently establishing ordered versus unordered type. For inline language and footnotes, the LibreOffice outputs contained relevant structural evidence while the Google Docs outputs did not. These examples show why a feature-level preservation comparison is more informative than a single destination-accessibility verdict.

\subsection{Confirmed losses}

The strongest negative findings are F06 and F11 under Google Docs. In F06, the source contains the run ``Buenos d\'ias'' with an OOXML \texttt{es-MX} override in an \texttt{en-US} document. PDF can represent this through a marked span with \texttt{/Lang}; the LibreOffice output does so. The Google Docs PDF is valid and tagged, but the destination extractor finds no inline span, raw PDF inspection finds no \texttt{es-MX} marker, and \texttt{pdftotext} confirms only the visible phrase. The outcome is \confirmedlost{} under the source contract.

In F11, the source contains a native footnote with an explicit reference and note body. The Google Docs PDF is valid and tagged and visibly contains the reference marker and note body, but raw inspection finds no \texttt{/Note}, \texttt{/Footnote}, or equivalent association marker, and the structural extractor returns no footnote. Because the PDF structure model can represent the feature and the other pipeline produced a comparable note structure, the result is \confirmedlost{} for the machine-verifiable footnote association under the source contract. This does not claim that the visible note text disappeared.

\begin{table*}[t]
\centering\scriptsize
\caption{Independent evidence for the two confirmed losses.}
\label{tab:losses}
\begin{tabular}{p{0.10\textwidth}p{0.20\textwidth}p{0.24\textwidth}p{0.28\textwidth}}
\hline
Case & Source evidence & Destination alternatives checked & Independent evidence / result \\
\hline
F06 / Google & OOXML \texttt{es-MX} run & PDF \texttt{/Lang}, \texttt{/ActualText}, raw marker scan & pypdf + PyMuPDF text; confirmed loss of inline language \\
F11 / Google & Native reference and footnote body & \texttt{/Note}, \texttt{/Footnote}, link association, raw objects & pypdf + PyMuPDF text; confirmed loss of machine-verifiable association \\
\hline
\end{tabular}
\end{table*}

\subsection{Measurement error and altered information}

F11 under LibreOffice is deliberately not described as preserved or lost. The output contains a \texttt{/Note} structure, but the extractor and independent text check cannot resolve the note-body association. This is a measurement error caused by insufficient destination evidence, not evidence of malformed output.

The three \altered{} cases change author-relevant information while retaining a related destination value: LibreOffice prefixes F02 alt text, Google Docs changes F05 from \texttt{en-US} to \texttt{en}, and Google Docs changes the F09 metadata title. Altered does not automatically mean unusable; it means that exact author-provided information was not retained unchanged.

\subsection{Integrated-document sanity check}

The secondary evaluation used three realistic two-page documents containing existing contract families: a community access report, an educational handout, and a policy note. The independent integrated source oracle passed all 21 property assertions. Across 42 property--pipeline observations, the protocol reported 16 verified preserved, 6 observed partial, 5 altered, 3 confirmed lost, 11 unresolved equivalence, and 1 measurement error. These observations are reported separately to show that the procedure remains usable when multiple properties coexist; they are not merged into the primary 26-case denominator and do not establish prevalence or a product ranking.



\section{Discussion}
\label{sec:discussion}

\subsection{What the benchmark demonstrates}

The primary result is methodological: accessibility preservation can be measured when source properties are explicitly known and destination representations can be inspected. The evidence-aware 26-case result contains 11 verified preserved cases, 3 observed partial cases, 3 alterations, 2 independently confirmed losses, 6 unresolved equivalence cases, and 1 measurement error. The categories do not collapse uncertainty into converter degradation: directly demonstrated changes are separated from cases where an alternate cross-format representation cannot yet be established.

This distinction matters because destination accessibility and source-information preservation answer different questions. A tagged PDF can contain headings, figures, or tables and still fail to preserve a particular source property. The F02 case is concrete: LibreOffice retained an alternative description but changed the author-provided string. Conversely, a missing destination structure can be a measurement problem rather than a converter loss, as shown by the LibreOffice F11 case. The comparison therefore complements, rather than replaces, destination-only accessibility evaluation.

\subsection{Converter-dependent behavior}

The same frozen DOCX source produced different historical V1 outcomes in 10 of 13 feature rows. The differences are property-specific rather than evidence for a universal winner. Google Docs retained F02 alt text exactly but confirmed loss of F06 inline language and F11 footnote structure in the tested pipeline. LibreOffice preserved explicit list-type evidence in F03 but altered F02 alt text. Such results support reporting the behavior of a named pipeline under a named configuration, not ranking products globally.

The Google findings should also be attributed carefully. The pipeline includes import and internal representation before PDF download. The data establish end-to-end behavior under the recorded workflow, not an isolated defect in Google's PDF renderer.

The integrated sanity check provides a limited external-validity signal: the same contracts and evidence statuses remained usable across three multi-feature documents. It is a secondary applicability demonstration, not a representative corpus or a new population estimate.

\subsection{Implications for authoring and conversion}

Standards and guidance already recognize preservation as a transformation concern. ATAG 2.0 includes a success criterion for preserving accessibility information during restructuring or recoding transformations when equivalent mechanisms exist, including text alternatives for non-text content \cite{atag20}. Revised Section 508 states that authoring tools should preserve information required for accessibility to the extent supported by the destination format \cite{section508}. EN 301 549 includes a corresponding requirement for preservation during transformations \cite{etsi301549}. A11yPreserve operationalizes this concern as a reproducible source/destination comparison for a constrained set of DOCX-to-PDF properties.

For document authors and publishing engineers, the practical implication is not that every export must be rejected or that one tool should always be selected. Rather, checking the destination should be complemented by checking whether critical author-supplied properties survived. The frozen benchmark suggests that this check can be automated for selected structures and that cases requiring human or improved parser review should remain explicitly marked.

\subsection{Contribution boundary}

A11yPreserve is a source-contract conformance protocol, benchmark, comparison tool, and empirical study. \texttt{a11ydiff} is the operational component, not the whole contribution. The reusable artifact is the combination of versioned contracts, independent source certificates, calibrated differential comparison, destination-parser triangulation, and explicit uncertainty statuses. The work does not attempt to build a general PDF checker, determine user preference, estimate population prevalence, or replace disabled-user evaluation. Its contribution is a concrete measurement chain from verified source information to destination structure and a conservative classification.



\section{Threats to Validity}
\label{sec:threats}

\subsection{Controlled corpus and feature units}

Atomic fixtures make source ground truth and causal interpretation clearer, but they are synthetic and small. Each feature family has one primary fixture, and feature families differ in complexity and accessibility relevance. Consequently, 26 cases are not IID samples and the outcome counts do not estimate real-world prevalence. The paper reports feature-level results first and gives the evidence-aware counts only as controlled-case summaries. Three integrated documents were evaluated as a separate sanity check, but they do not make the corpus representative; justified variants and broader source corpora remain future work.

\subsection{Converter and format coverage}

The experiment evaluates DOCX-to-PDF only, using two real pipelines. It does not characterize Microsoft Word, ODT, HTML, EPUB, or other conversion routes. Two independent pipelines are sufficient to demonstrate that behavior can vary across pipelines, but not to characterize the converter ecosystem. Microsoft Word was deferred because the environment did not yield a completed primary workflow; it is excluded rather than treated as a semantic result.

\subsection{Cloud reproducibility}

Google Docs is a cloud service whose backend version cannot necessarily be pinned. The study records the date, operating system, browser, observed renderer, workflow, source hashes, output hashes, and validation evidence. A rerun can nevertheless differ after a service update. The correct interpretation is therefore a dated Google Docs DOCX-to-PDF pipeline condition, not a timeless property of Google Docs.

\subsection{Cross-format representation}

Source and destination formats express accessibility information differently. The study uses predefined feature contracts and machine-inspectable mappings, such as Word Heading 2 to PDF \texttt{/H2}, a Word run-language override to PDF \texttt{/Lang}, or a native equation to PDF \texttt{/Formula}. Some mappings are incomplete or parser-sensitive. F07, F08, F10, F12, F14, and especially F11 retain explicit evidence limitations. The measurement-error category prevents such limitations from being silently converted into loss or preservation.

\subsection{Parser and validator limitations}

The primary destination extractor is a structured parser, supplemented by raw object inspection and Poppler text extraction. Independent checks reduce but do not eliminate parser error. In particular, absence from one extracted encoding is not automatically proof of semantic absence when equivalent cross-format representations are possible. Conventional accessibility validators were not available as a systematic dataset and are not treated as ground truth. The study therefore does not claim complete PDF accessibility, that validators are defective because they miss preservation changes, or that \texttt{a11ydiff} measures complete accessibility.

\subsection{User impact}

No disabled-user study, screen-reader task study, or perceived-usability study was performed. The results show whether encoded source information was found in the destination representation under the contracts; they do not measure task completion, screen-reader usability, perceived accessibility, or severity for disabled users. A future user-centered study could connect preservation changes to task impact.

\subsection{Interpretation and generalization}

An \altered{} description may remain usable, an \observedpartial{} structure may still provide substantial accessibility, and an \unresolved{} case may ultimately prove equivalent under a richer representation. Conversely, this study does not determine the practical severity of any individual change. The classifications are evidence-centered preservation labels, not a weighted accessibility score. Claims are limited to the frozen atomic fixtures, the three secondary integrated documents, and recorded pipeline conditions: they do not establish that conversion generally destroys accessibility, that Google Docs is generally worse or LibreOffice generally better, that findings generalize to all DOCX documents or software versions, that altered information necessarily causes user harm, or that PDF export alone caused every observed change.



\section{Conclusion}
\label{sec:conclusion}

A11yPreserve demonstrates a practical source-grounded way to measure accessibility-information preservation across document conversion. The study froze 13 controlled DOCX accessibility fixtures, verified their source representations, converted them through LibreOffice and the Google Docs DOCX-to-PDF pipeline, inspected destination PDF structure, and compared the paired representations with retained evidence and hashes.

Across 26 real outputs, the evidence-aware benchmark recorded 11 verified preserved, 3 observed partial, 3 altered, 2 confirmed lost, 6 unresolved equivalence, and 1 measurement-error case. The two confirmed losses were the inline-language and footnote properties in the Google Docs pipeline. A separate three-document sanity check showed that the same protocol remained usable when multiple existing feature contracts coexisted. The study also found that a tagged destination can retain accessibility structure while changing or failing to preserve particular source information. This is why destination accessibility and source accessibility preservation should be reported as related but distinct questions.

The conclusions are intentionally bounded. The corpus is controlled and not representative, only two DOCX-to-PDF pipelines were evaluated, Google Docs is cloud-versioned, PDF representation equivalence can be parser-sensitive, and no disabled-user study was performed. Within those limits, the evidence supports the feasibility of source-grounded preservation measurement and shows that preservation outcomes differed across tested features and the two evaluated pipelines. It does not support claims about general accessibility, a global product ranking, all DOCX documents, all software versions, universal user harm, or PDF-export-only causation. Future work can extend the corpus with justified variants, compare additional converters, study validator detection, and connect structural changes to user tasks without changing the central question: did known accessibility information survive conversion?




\bibliographystyle{plain}
\bibliography{references}

\end{document}

